\clearpage\corrigeDevoir{5}
\begin{blocLivre}{dscor-5-1}\begin{corrigeBac}{6}{Son et distance}
\demarche{Ne pas confondre retard de propagation et période du signal.}
\schema[corrige]{retard-son}
\begin{questions}
\item $v=1,02/0,00300=340$ m/s.\pts{2 pt}
\item $T=10,0/5=2,00$ ms ; $f=500$ Hz.\pts{2 pt}
\item Hauteur : liée à la fréquence ; timbre : lié à la forme ; niveau : lié à l’intensité sur l’échelle des décibels.\pts{2 pt}
\end{questions}
\end{corrigeBac}\end{blocLivre}
\begin{blocLivre}{dscor-5-2}\begin{corrigeBac}{7}{Spectre et réfraction}
\demarche{Appliquer chaque loi à la bonne interface et mesurer les angles à la normale.}
\schema[corrige]{refraction}
\begin{questions}
\item Non ; $520\unit{nm}=\qty{5,20e-7}{m}$.\pts{2 pt}
\item $\sin r=\sin30^\circ/1,50=1/3$, donc $r\simeq19,5^\circ$.\pts{2 pt}
\item $i'=i=30,0^\circ$, mesuré à la normale.\pts{1 pt}
\item L’indice dépend de la longueur d’onde ; les radiations sont réfractées différemment.\pts{2 pt}
\end{questions}
\end{corrigeBac}\end{blocLivre}
\begin{blocLivre}{dscor-5-3}\begin{corrigeBac}{7}{Une image nette}
\demarche{Construire avant de lire, puis vérifier les tailles par Thalès.}
\schema[corrige]{lentille-3f}
\begin{questions}
\item Le rayon central et le rayon parallèle réfracté par le foyer image se coupent au point image.\pts{2 pt}
\item Image à 15 cm, de taille 1,0 cm et renversée ; $\gamma=-0,50$.\pts{2 pt}
\item À 15 cm après la lentille, dans le plan de l’image.\pts{1 pt}
\item Lentille : ensemble optique de l’œil ; écran : rétine ; ajouter un diaphragme pour la pupille.\pts{2 pt}
\end{questions}
\end{corrigeBac}\end{blocLivre}

